\topic{2 continued}{Why pressure and velocity are staggered}
In the preceding grid, a cell center is where pressure is stored. The shared
face is the boundary between two adjacent cells, where the horizontal velocity
is stored. Along one horizontal row, let $x_i$ be center $i$'s position and
$p_i$ its pressure sample. The midpoint $x_{i+1/2}=(x_i+x_{i+1})/2$ is the
position of the face between centers $i$ and $i+1$.
\lessonfigure{checkerboard_pressure}{A synthetic pressure diagnostic with h=0.25. (a) Pressure samples at equally spaced cell centers alternate between +1 and -1. (b) Blue circles use the two samples on either side of a cell center; red crosses use the two adjacent samples around their midpoint. The blue result is zero while the red results are +8 or -8. The legend identifies these two difference formulas and their evaluation positions; it does not show two flow solutions.}{checkerboard-pressure}
Panel (b) compares \emph{where} a pressure derivative is evaluated and
\emph{which pair} of stored values is used. Blue circles refer to a centered
derivative at a cell center; red crosses refer to a derivative at the
midpoint between adjacent centers. Their definitions are expanded next.
\newpage
\topic{2 continued}{Read the two pressure-gradient formulas}
Consider a synthetic interior grid patch with $p_i=(-1)^i$. The index $i$
counts pressure samples along $x$; the spacing is $h=0.25$. Neighboring
pressures alternate between $+1$ and $-1$.

If velocity and pressure were both stored at the same centers and the pressure
gradient were approximated by a naive central difference there, it would give
\[
\left(\frac{\partial p}{\partial x}\right)_i
\simeq\frac{p_{i+1}-p_{i-1}}{2h}=0,
\]
because $p_{i+1}=p_{i-1}$. This alternating pressure mode is invisible to that
gradient: the resulting momentum equation cannot distinguish it from zero
pressure gradient in the interior.

On the MAC grid, horizontal momentum is evaluated at the shared face between
adjacent pressure centers. Its pressure gradient is
\[
\left(\frac{\partial p}{\partial x}\right)_{i+1/2}
\simeq\frac{p_{i+1}-p_i}{h}=\pm\frac{2}{h}=\pm8.
\]
The face momentum equation responds to the adjacent pressure difference.
This is one reason for staggering the unknowns. In two dimensions an alternating
pattern in both directions is called a pressure checkerboard.

The figure is a discrete diagnostic pattern, not a computed cavity solution.
Properly stabilized collocated methods can also work; they require their own
face-flux interpolation treatment. The present course implementation uses the
staggered arrangement derived here.
